% A Casimir-locked seesaw for electroweak symmetry breaking
% Prepared for submission to Physical Review D.

\documentclass[%
 aps,
 prd,
 reprint,
 nofootinbib,
 superscriptaddress,
 amsmath,
 amssymb,
]{revtex4-2}

\usepackage{graphicx}
\usepackage{hyperref}
\usepackage{xcolor}
\hypersetup{colorlinks=true,linkcolor=blue,citecolor=blue,urlcolor=blue}

% ---------------------------------------------------------------
\begin{document}

\title{A Casimir-locked seesaw for electroweak symmetry breaking}

\author{[Author]}
\affiliation{[Affiliation]}

\date{\today}

\begin{abstract}
We present a minimal algebraic ansatz in which the electroweak gauge-boson and bare-Higgs mass spectrum is generated by a pair of $2\times 2$ mass-squared matrices labelled by the two lowest non-trivial irreducible representations of $\mathrm{SU}(2)_W$. Each matrix has the Casimir-locked form $M_R^2 = m_0^2\,\bigl[\,[0,\sqrt{C_2(R)}\,],[\sqrt{C_2(R)},-C_2(R)\,]\,\bigr]$ and shares a single seed scale $m_0$. Anchoring the doublet upper eigenvalue to $M_W^2$ fixes $m_0=106.5702$~GeV; the same one-parameter ansatz then predicts $M_Z=91.1814$~GeV (PDG $91.1876$, deviation $-0.0068\%$) and a bare Higgs mass $m_{h,\mathrm{bare}}=122.3809$~GeV. A rank-1, eigenstate-aligned spurion $\Delta=(26.4189\,\mathrm{GeV})^2$, naturally generated by the one-loop top Coleman--Weinberg shift at cutoff $m_0$, lifts $m_{h,\mathrm{bare}}$ to $m_h=125.20$~GeV. The construction is purely electroweak: it predicts a four-eigenvalue spectrum from one seed scale plus one trace-shifting spurion, and is silent on fermion masses, gauge couplings, and the unbroken $\mathrm{U}(1)_{\mathrm{EM}}$ subgroup. The seed structure, $b=1$ coefficient, and the spurion mechanism are postulates rather than derived from a unique UV completion.
\end{abstract}

\maketitle

% ---------------------------------------------------------------
\section{Introduction}

The Standard Model leaves the Yukawa couplings and the bare Higgs mass-squared as inputs. The first set is reverse-engineered from twelve measured fermion masses and mixing angles; the second is the unstable scale whose smallness against any UV cutoff is the hierarchy problem. Any candidate for a structure underneath the SM should reduce the number of independent dimensionful inputs, predict at least one observable that is currently fitted, and do so without introducing more free parameters than it removes.

We pursue here a minimal answer to a narrower question --- the \emph{wherefore} problem for electroweak symmetry breaking. We postulate that the bare electroweak mass-squared spectrum is the diagonalisation of two $2\times 2$ matrices, one per chosen $\mathrm{SU}(2)_W$ representation, of a single Casimir-locked form sharing one seed scale $m_0$. We then ask which representations are selected, what fixes $m_0$ numerically, and what algebraic object plays the role of the small shift that brings the bare Higgs onto its physical pole. The answers are: the doublet and the triplet uniquely (both as the two lowest non-trivial $\mathrm{SU}(2)$ irreps and as the unique numerical winner among ten admissible pairs); $m_0$ is anchored by $M_W$ and predicts $M_Z$ to $7\times10^{-5}$; the operative spurion is rank-1 along the tachyonic eigenvector, of the size produced by the top-loop Coleman--Weinberg shift at cutoff $m_0$.

The construction is deliberately restricted to the electroweak boson sector. The fermion sector --- Yukawa hierarchy, Koide-type lepton relations, CKM and PMNS mixing --- is \emph{not} addressed in this paper. A previous version included a Hermitian $\mathbb{Z}_3$ circulant for the charged leptons and an attempted extension to the $(s,c,b)$ down-quark tuple; those constructions are an empirical parametrisation of Koide-type relations and have no derivable connection to the Casimir-locked seed of Eq.~(\ref{eq:seed}). They have been excised. The natural next step --- a Lagrangian completion that ties a fermion-mass spurion to the rank-1 boson spurion of Sec.~\ref{sec:ewfit} --- is left as future work.

The ansatz makes contact with two pre-existing lines of work. The SO(32) scalar classification of Rivero~\cite{Rivero2024SO32} supplies a parent group in which the doublet--triplet pair can be embedded via a standard Slansky chain; the sBootstrap framework~\cite{Rivero2005sBoot} supplies a class of composite-Higgs UV completions in which the required spurion is generic in size and sign. The closest mechanistic analogues in the bosonic-seesaw literature --- Calmet--Oliver~\cite{CalmetOliver2007}, Kim~\cite{Kim2005}, Chivukula--Dobrescu--Georgi--Hill~\cite{Chivukula1999}, and Haba \emph{et al.}~\cite{Haba2016} --- are reviewed in Sec.~\ref{sec:corpus_loop}. The relations are structural rather than dynamical: the ansatz is not derived from any of these starting points.

The paper is organised as follows. Section~\ref{sec:seed} introduces the Casimir-locked seed and derives its general form. Section~\ref{sec:repsel} explains why the doublet--triplet pair is selected. Section~\ref{sec:ewfit} reports the electroweak fit, the size and sign of the spurion $\Delta$, and its mechanism. Section~\ref{sec:corpus} makes the embedding into the existing corpus precise. Section~\ref{sec:caveats} lists open issues; Sec.~\ref{sec:conclusion} concludes.

% ---------------------------------------------------------------
\section{The Casimir-locked seed}\label{sec:seed}

\subsection{Definition}

For each $\mathrm{SU}(2)_W$ representation $R$ of dimension $n$ we define a real symmetric $2\times 2$ mass-squared matrix
\begin{equation}\label{eq:seed}
M_R^2 \;=\; m_0^2 \begin{pmatrix} 0 & a\,\sqrt{C_2(R)} \\ a\,\sqrt{C_2(R)} & -\,b\,C_2(R)\end{pmatrix},\qquad C_2(n)=\frac{n^2-1}{4}.
\end{equation}
The quadratic Casimir is normalised so that $T^a T^a = C_2(R)\,\mathbf{1}_R$ with Hermitian generators satisfying $[T^a,T^b]=i\epsilon^{abc}T^c$; the doublet (fundamental) has $T^a=\sigma^a/2$ and $C_2(2)=3/4$, while the adjoint triplet has $C_2(3)=2$. The metric is mostly-plus; a positive eigenvalue of $M_R^2$ is a physical mass-squared, a negative one is a tachyonic slot that requires a separate condensation mechanism. The seed scale $m_0$ has units of GeV and is the only dimensional parameter in the construction. The coefficients $a$ and $b$ are dimensionless real numbers; throughout this paper we take $a=b=1$, and discuss this choice below.

\subsection{The three conditions and the most general form}

Three conditions specify Eq.~(\ref{eq:seed}) up to the coefficients. (i)~$\det M_R^2 < 0$, so that one eigenvalue is tachyonic and the seesaw has a candidate Higgs slot. (ii)~The trace $\mathrm{Tr}\,M_R^2$ is a polynomial in $C_2(R)$ with leading order $C_2(R)$, and no constant term. (iii)~A discrete chiral or gauge symmetry protects one diagonal entry, conventionally the $(1,1)$ slot, against acquiring a bilinear mass. Under (iii) one diagonal entry is zero. Under (ii) plus the additional postulate that the polynomial has no constant term, the other diagonal entry is $-b\,C_2(R)\,m_0^2$. Condition (i) is then automatic provided the off-diagonal entry is non-zero. The minimal-power-of-Casimir choice $m_{12}=a\sqrt{C_2(R)}\,m_0^2$ is the unique odd-power form consistent with the group-theoretic content; higher odd powers ($C_2^{3/2}$, etc.) are excluded by minimality.

\subsection{Status of $a=b=1$}

The coefficient $a$ can be absorbed into a rescaling of $m_0$ by canonically normalising the kinetic term of the protected $(1,1)$ field; setting $a=1$ is therefore a convention, not a prediction. The coefficient $b$ has no analogous reparameterisation. None of the natural group-theoretic candidates we examined --- a Wess--Zumino-style covariant-derivative argument, a custodial-$\mathrm{SU}(2)$ Ward identity, a trace condition on $M_R^2$, or a $(1,1)$ kinetic-normalisation --- rigorously fixes $b$. We therefore record $b=1$ as an independent postulate, the simplest non-trivial integer Wilson coefficient in the polynomial expansion of $m_{22}$ in $C_2$. The ``Casimir-locked'' label refers to the \emph{shape} of Eq.~(\ref{eq:seed}) --- both entries scale with the same $C_2(R)$ with no $R$-dependent dimensionless coefficients --- not to a derivation of $b$.

\subsection{Eigenvalues}

With $t=\sqrt{C_2(R)}$ and $a=b=1$, the eigenvalues of Eq.~(\ref{eq:seed}) in units of $m_0^2$ are
\begin{equation}\label{eq:eig}
\lambda_{\pm}(R) \;=\; \tfrac{1}{2}\Bigl(-C_2(R) \pm \sqrt{C_2(R)^2 + 4 C_2(R)}\Bigr).
\end{equation}
For the doublet and triplet, this gives in closed form
\begin{equation}\label{eq:eig23}
\begin{aligned}
\lambda_{\pm}(2) &= \tfrac{\pm\sqrt{57}-3}{8} \in \{+0.5687,\;-1.3187\},\\
\lambda_{\pm}(3) &= \pm\sqrt{3}-1 \in \{+0.7321,\;-2.7321\}.
\end{aligned}
\end{equation}
These four real numbers carry the entire electroweak content of the ansatz once $m_0$ is fixed.

% ---------------------------------------------------------------
\section{Representation selection}\label{sec:repsel}

\subsection{Rule~A: two lowest non-trivial $\mathrm{SU}(2)$ irreps}

The simplest parameter-free, purely group-theoretic rule that uniquely produces the pair $\{R_1,R_2\}=\{2,3\}$ is to take the two lowest non-trivial irreducible representations of $\mathrm{SU}(2)$, i.e.\ $n=2$ and $n=3$. This rule has no free parameter and does not invoke any phenomenological input. We refer to it as Rule~A. Four other rules were considered (phenomenological matter-content, minimisation of $|C_2(R_1)+C_2(R_2)-k|$ for some integer $k$, custodial $\mathrm{SU}(2)_L\times\mathrm{SU}(2)_R/\mathbb{Z}_2$ decomposition, and a two-mode KK truncation of a 5D $\mathrm{SU}(2)$ theory); each either is circular, requires extra parameters, or numerically fails.

\subsection{Ten-pair enumeration}\label{sec:scan}

There are ten unordered pairs of $\mathrm{SU}(2)$ irreps with dimension $n\in\{2,3,4,5,6\}$ (equivalently $C_2<12$). The scoring rule is fixed once: for each pair we anchor the best available eigenvalue to $M_W^2$ (this fits $m_0$), and we report
\begin{equation}\label{eq:scoring}
\mathcal{E}(R_1,R_2)\;:=\;\max_{X\in\{M_Z,\,m_{h,\mathrm{bare}}\}}\;\bigl|\,X_{\mathrm{pred}}/X_{\mathrm{target}}-1\,\bigr|,
\end{equation}
i.e.\ the maximum fractional deviation of the two non-anchored target masses $\{M_Z,m_{h,\mathrm{bare}}\}$ from their reference values (PDG $M_Z=91.1876$~GeV and $m_{h,\mathrm{bare}}=122.4$~GeV). Each pair has four eigenvalues; the two unused eigenvalues do not enter $\mathcal{E}$ and are not constrained by the scan. We made this choice because $\{M_W,m_h,v/\sqrt{2}\}$ are model inputs (cf.\ Sec.~\ref{sec:ewfit}) and so only $M_Z$ and $m_{h,\mathrm{bare}}$ are genuine predictions to score against. The Higgs-pole $m_h$ requires the spurion $\Delta$, which is a separate fit and so is excluded from $\mathcal{E}$.

\begin{table}[t]
\caption{Ten-pair enumeration scored by $\mathcal{E}$ of Eq.~(\ref{eq:scoring}). For each pair, $m_0$ is fitted to anchor the best-aligned eigenvalue to $M_W$, and $\mathcal{E}$ is the maximum fractional deviation across $\{M_Z, m_{h,\mathrm{bare}}\}$.}
\label{tab:pairs}
\begin{ruledtabular}
\begin{tabular}{cccc}
pair & $(C_2^{(1)},C_2^{(2)})$ & best $m_0$ (GeV) & $\mathcal{E}$ \\
\hline
$(2,3)$ & $(0.75,\,2.00)$ & $106.570$ & $\mathbf{0.02\%}$ \\
$(2,4)$ & $(0.75,\,3.75)$ & $106.570$ & $5.86\%$ \\
$(2,5)$ & $(0.75,\,6.00)$ & $106.570$ & $9.19\%$ \\
$(2,6)$ & $(0.75,\,8.75)$ & $106.570$ & $11.25\%$ \\
$(3,4)$ & $(2.00,\,3.75)$ & $\phantom{0}88.727$ & $19.82\%$ \\
$(3,5)$ & $(2.00,\,6.00)$ & $\phantom{0}86.017$ & $19.29\%$ \\
$(3,6)$ & $(2.00,\,8.75)$ & $\phantom{0}84.428$ & $20.78\%$ \\
$(4,5)$ & $(3.75,\,6.00)$ & $\phantom{0}86.017$ & $50.24\%$ \\
$(4,6)$ & $(3.75,\,8.75)$ & $\phantom{0}84.428$ & $47.46\%$ \\
$(5,6)$ & $(6.00,\,8.75)$ & $\phantom{0}84.428$ & $80.83\%$ \\
\end{tabular}
\end{ruledtabular}
\end{table}

The pair $\{2,3\}$ is the unique numerical winner by a factor of $\sim 300$ in maximum fractional error. The dominance has two independent algebraic sources. First, only $n=2$ delivers $\sqrt{|\lambda_-/\lambda_+|}\approx 1.523$, the ratio that maps $M_W=80.369$~GeV to $m_{h,\mathrm{bare}}=122.4$~GeV; higher $n$ overshoots monotonically. Second, once $R_1=2$ is fixed, the requirement $\sqrt{\lambda_+(R_2)/\lambda_+(2)}=M_Z/M_W=1.1346$ is satisfied only by $\lambda_+(R_2)=\sqrt{3}-1$, which is the triplet alone.

% ---------------------------------------------------------------
\section{Electroweak fit and the spurion $\Delta$}\label{sec:ewfit}

\subsection{Fit~B: doublet upper eigenvalue $=M_W^2$}\label{sec:fitB}

We identify the doublet $\lambda_+(2)\,m_0^2$ with the physical $W$-boson mass-squared and the doublet $\lambda_-(2)\,m_0^2$ with the negative of the bare Higgs mass-squared:
\begin{equation}\label{eq:fitB}
\lambda_+(2)\,m_0^2 \;=\; M_W^2,\qquad |\lambda_-(2)|\,m_0^2 \;=\; m_{h,\mathrm{bare}}^2.
\end{equation}
Imposing the first relation fixes
\begin{equation}\label{eq:m0}
m_0 \;=\; \frac{M_W}{\sqrt{(\sqrt{57}-3)/8}} \;=\; 106.5702\;\mathrm{GeV}.
\end{equation}
The triplet eigenvalues then predict, with no further free parameter,
\begin{equation}\label{eq:MZ}
\begin{aligned}
M_Z &= \sqrt{\sqrt{3}-1}\,m_0 = 91.1814\;\mathrm{GeV},\\
|\lambda_-(3)|^{1/2}\,m_0 &= 176.149\;\mathrm{GeV},
\end{aligned}
\end{equation}
and the doublet lower eigenvalue gives $m_{h,\mathrm{bare}}=122.3809$~GeV.

\begin{table}[t]
\caption{Verification of Fit~B at four-decimal precision (\texttt{mpmath}, $\texttt{mp.dps}=30$).}
\label{tab:fit}
\begin{ruledtabular}
\begin{tabular}{lcccl}
Quantity & Predicted (GeV) & Target (GeV) & Deviation & Role \\
\hline
$m_0$           & $106.5702$ & $106.5700$ (seed) & $+0.0002\%$ & anchored by $M_W$ \\
$M_W$           & $\phantom{0}80.3690$ & $\phantom{0}80.3690$ (PDG) & $0\%$ & anchor \\
$M_Z$           & $\phantom{0}91.1814$ & $\phantom{0}91.1876$ (PDG) & $-0.0068\%$ & prediction \\
$m_{h,\mathrm{bare}}$ & $122.3809$ & $122.4$ (seed) & $-0.0156\%$ & prediction \\
$m_h$           & $125.2000$ & $125.20$ (PDG) & $\approx 0$ & anchored by $\Delta$ \\
$v/\sqrt{2}$    & $174.1036$ & $174.1036$ ($G_F$) & $0$ & external input \\
\end{tabular}
\end{ruledtabular}
\end{table}

Honest accounting of inputs vs.\ outputs: $m_0$ is fixed by anchoring to $M_W$, $\Delta$ is fixed by anchoring to $m_h$, and $v/\sqrt{2}$ is taken from $G_F$. With two model parameters $(m_0,\Delta)$ and three external anchors $\{M_W,m_h,G_F\}$, the genuine \emph{predictions} of the construction are the two remaining table rows: $M_Z=91.1814$~GeV ($-0.0068\%$ vs.\ PDG, i.e.\ $6.2$~MeV) and $m_{h,\mathrm{bare}}=122.3809$~GeV ($-0.0156\%$ vs.\ the seed-quoted reference value, which is itself defined only modulo the choice of bare-vs.-pole convention). Both predictions are at the $10^{-4}$ level. The construction is therefore better described as a \emph{one-parameter fit to $M_W$ that simultaneously reproduces $M_Z$} (and a separate one-parameter fit of $\Delta$ to $m_h$), not as a four-observable fit with two parameters.

\subsection{The spurion: rank-1, not trace-preserving}\label{sec:spurion}

The shift required to bring $m_{h,\mathrm{bare}}=122.3809$~GeV to $m_{h,\mathrm{phys}}=125.20$~GeV is
\begin{equation}\label{eq:Delta}
\Delta \;:=\; m_{h,\mathrm{phys}}^2 - m_{h,\mathrm{bare}}^2 \;=\; (26.4189\,\mathrm{GeV})^2.
\end{equation}
An earlier description of $\Delta$ as a ``trace-preserving spurion'' is technically incorrect. A literal trace-preserving perturbation $\delta\lambda_+ + \delta\lambda_-=0$ with $\delta\lambda_-=-(26.42\,\mathrm{GeV})^2$ would shift $\delta\lambda_+=+(26.42\,\mathrm{GeV})^2$ and move $M_W$ from $80.369$~GeV to $84.600$~GeV --- a $4.23$~GeV mismatch, $\sim 325\sigma$ outside the PDG value. A determinant-preserving shift fails by $\sim 1.9$~GeV.

The only shift consistent with $m_h\!:\!122.38\to 125.20$~GeV and $M_W$ left at $80.369$~GeV to PDG precision is a \emph{rank-1, eigenstate-aligned} spurion proportional to the projector $P_-=|v_-\rangle\langle v_-|$ onto the tachyonic eigenvector of Eq.~(\ref{eq:seed}). In the $(a=b=1)$ doublet basis,
\begin{equation}\label{eq:P}
P_- \;=\; \begin{pmatrix} \tfrac{1}{2}-\tfrac{\sqrt{57}}{38} & -\tfrac{2\sqrt{19}}{19} \\ -\tfrac{2\sqrt{19}}{19} & \tfrac{1}{2}+\tfrac{\sqrt{57}}{38}\end{pmatrix},\qquad \mathrm{Tr}\,P_-=1,\;\det P_-=0,
\end{equation}
and the operative perturbation is $\delta M_2^2 = -\Delta\,P_-$ with $\Delta>0$. This shift is spectator-preserving: it leaves $\lambda_+=M_W^2$ untouched and lowers $\lambda_-$ by exactly $\Delta$, lifting $m_h$ by $\Delta$ in mass-squared.

\subsection{Mechanism for $\Delta$}\label{sec:mech}

Five candidate mechanisms were tested against $\Delta=(26.42\,\mathrm{GeV})^2$. Two land within an order of magnitude with at most one TeV-scale auxiliary; the remaining three miss by factors $27$ to $1500$ at their natural scales.

\paragraph{M4 --- one-loop top Coleman--Weinberg at cutoff $\Lambda=m_0$.} The quadratically-divergent top-loop contribution, regulated by a hard cutoff at the seed scale $m_0$, evaluates to
\begin{equation}\label{eq:M4}
\Delta_{M4,\,\mathrm{quad}} \;=\; -\frac{3y_t^2}{8\pi^2}\bigl(m_0^2 - m_t^2\bigr) \;=\; (26.2766\,\mathrm{GeV})^2,
\end{equation}
using $y_t=\sqrt{2}\,m_t/v_{\mathrm{SM}}=0.99196$, $m_0=106.57$~GeV, and $m_t=172.7$~GeV. The result is $0.989$ of the required $(26.4189\,\mathrm{GeV})^2$ --- agreement at the $1\%$ level with no auxiliary scale. The seed scale fixed by the doublet anchor in Sec.~\ref{sec:fitB} turns out to be exactly the cutoff at which the top loop produces $\Delta$ of the right size. The result is regulator-dependent (the quadratic divergence is not defined in $\overline{\mathrm{MS}}$); the agreement may be coincidence, or it may indicate that $m_0$ is the physical matching scale of the construction. The log piece of the same Coleman--Weinberg expression gives $\Delta_{M4,\,\mathrm{log}}=(32.82\,\mathrm{GeV})^2$, $1.54\times$ the target --- within an order of magnitude but a poorer fit.

\paragraph{M3 --- dim-6 $|H|^4$ EFT alternative.} A $c_6(H^\dagger H)^3/\Lambda_d^2$ operator generated at a UV scale $\Lambda_d$ shifts $m_h^2$ by $\Delta_{M3}\sim c_6\,(v/\sqrt{2})^4/\Lambda_d^2$. For $c_6=1$, the required $\Delta=(26.42\,\mathrm{GeV})^2$ corresponds to $\Lambda_d=1.15$~TeV; for $c_6=4$ (the standard normalisation of the $(H^\dagger H)^3$ operator), $\Lambda_d=2.30$~TeV. Both lie in the TeV range and are consistent with current EW precision and direct-search bounds. M3 is the natural EFT statement: any UV completion that couples to the Higgs as a dim-6 operator at $\Lambda_d\sim 1$--$2$~TeV generates $\Delta$ of the right size.

The remaining three mechanisms fall short. M1 (tree-level SUSY soft mass) misses by $\sim 1500\times$ at natural SUSY scales and requires anomaly-mediated tuning at $\Lambda_{\mathrm{SUSY}}\sim 10$~TeV. M2 (compositeness threshold at $\Lambda_c\sim 1$~TeV) misses by $\sim 27\times$. M5 (trace-anomaly piece $|\beta_\lambda|\Lambda^2$) is numerically degenerate with M4 and not independent.

% ---------------------------------------------------------------
\section{Connection to the existing corpus}\label{sec:corpus}

\subsection{SO(32) embedding}

Rivero~\cite{Rivero2024SO32} proposes an SO(32) classification of generic supersymmetric scalars using a \emph{flavour} $\mathrm{SU}(5)_F$ chain, not $\mathrm{SU}(2)_W$. We therefore cannot extract the $\{2,3\}$ pair from the conventions of Ref.~\cite{Rivero2024SO32} directly. A standard Slansky~\cite{Slansky1981} decomposition, however, places the seesaw pair inside SO(32) via the maximal chain
\begin{equation}\label{eq:so32}
\mathrm{SO}(32)\;\supset\;\mathrm{SO}(28)\times\mathrm{SO}(4)\;\supset\;\mathrm{SO}(28)\times\mathrm{SU}(2)_W^{\mathrm{diag}},
\end{equation}
with the adjoint $\mathbf{496}$ branching as $(\mathbf{378},1,1)+(\mathbf{1},3,1)+(\mathbf{1},1,3)+(\mathbf{28},2,2)$ and the bidoublet $(2,2)$ decomposing under the diagonal $\mathrm{SU}(2)_W$ as $\mathbf{1}+\mathbf{3}$. The triplet of $\mathrm{SU}(2)_W$ is contained in both $(\mathbf{1},3,1)+(\mathbf{1},1,3)$ and in the bidoublet piece; the doublet of $\mathrm{SU}(2)_W$ sits in the SO(32) spinor $\mathbf{32768}$. The chain~(\ref{eq:so32}) shows that $\{2,3\}$ fits naturally inside an SO(32) parent --- but it is not uniquely selected by the embedding, because Eq.~(\ref{eq:so32}) also produces $\mathrm{SU}(2)_W$ singlets and the choice of which subsector to retain is external.

\subsection{sBootstrap and the size of $\Delta$}

Rivero's sBootstrap proposal~\cite{Rivero2005sBoot} places the bosonic superpartners of the SM fermions in composite quark--diquark states, naturally producing soft masses of order $v/\sqrt{2}$. The result $\Delta=(26.4189\,\mathrm{GeV})^2=(v/\sqrt{2})^2/43$ is \emph{smaller} than a natural tree-level sBootstrap soft mass by a factor $5$--$40$, but lands at the right size once a one-loop top-loop suppression is applied:
\begin{equation}\label{eq:sB}
\Delta_{\mathrm{sB}}\;=\;\frac{3y_t^2}{8\pi^2}\,\Bigl(\tfrac{v}{\sqrt{2}}\Bigr)^2\;c_{\mathrm{sB}},
\end{equation}
with $c_{\mathrm{sB}}\simeq 0.61$. Equation~(\ref{eq:sB}) is the same physics as the top-loop mechanism M4 of Sec.~\ref{sec:mech}, written in sBootstrap variables; we do not claim that Ref.~\cite{Rivero2005sBoot} derived Eq.~(\ref{eq:sB}). The compatibility is in size and sign.

\subsection{Bosonic-seesaw analogues}\label{sec:corpus_loop}

The $2\times 2$ scalar mass-squared structure of Eq.~(\ref{eq:seed}) is closely related to several pre-existing ``bosonic seesaw'' constructions. Calmet--Oliver~\cite{CalmetOliver2007} use a $2\times 2$ mass-squared matrix in a two-Higgs-doublet sector whose tachyonic eigenvalue drives EWSB; the \emph{shape} of their matrix is the closest published analogue of Eq.~(\ref{eq:seed}) restricted to the doublet sector, but their construction has neither the Casimir locking nor the triplet sector. Kim~\cite{Kim2005} coins the term ``bosonic see-saw'' for a scalar $2\times 2$ matrix whose negative eigenvalue is the SM Higgs, motivated by SUSY breaking; the algebraic skeleton is the same as our doublet but the seed scale is set by SUSY mediation rather than by group-theory invariants. Chivukula--Dobrescu--Georgi--Hill~\cite{Chivukula1999} use a structurally identical $2\times 2$ matrix on the fermion side (the top-quark seesaw), supplying a direct fermion-side analogue. Haba \emph{et al.}~\cite{Haba2016} generate the negative mass-squared of a bosonic seesaw \emph{radiatively} via Coleman--Weinberg dynamics; this is the closest mechanistic precursor of the M4 top-loop spurion of Sec.~\ref{sec:mech}, although Ref.~\cite{Haba2016} does not consider a Casimir-locked algebraic skeleton.

% ---------------------------------------------------------------
\section{Caveats and open issues}\label{sec:caveats}

We retain the adversarial findings as honest disclosure. The construction has the following load-bearing weaknesses.

\paragraph{(a) The seed structure is engineered, not derived.} The matrix Eq.~(\ref{eq:seed}) has the $(1,1)$ entry zero, the $(2,2)$ entry equal to $-C_2(R)\,m_0^2$, and the off-diagonal equal to $\sqrt{C_2(R)}\,m_0^2$. There is no symmetry argument that uniquely forces this assignment; the Casimir appears at three positions with three different powers $\{0,\tfrac12,1\}$. Replacing the $(1,1)$ zero with $\epsilon m_0^2$ for any $\epsilon=O(0.05)$ breaks the eigenvalue ratio that delivers $M_W/m_0$ and $m_{h,\mathrm{bare}}/m_0$. The model has no structural robustness against the most natural perturbations of its own ansatz. An exploratory analysis of candidate renormalisable Lagrangians finds the following partial picture. The $(1,1)=0$ slot is naturally protected by a $\mathbb{Z}_2$ or by chiral protection of an elementary scalar in a partial-compositeness setting~\cite{Caracciolo2013PC}. The $(2,2)=-C_2(R)\,m_0^2$ scaling is reproduced automatically by $\Phi_R^\dagger T^a T^a\Phi_R$, either as a renormalisable mass operator or as the composite-mass term arising from gauging $\mathrm{SU}(2)_W$ inside the global symmetry of a confining hypercolor sector. The $\sqrt{C_2(R)}$ off-diagonal, however, is the structural obstruction: it is not produced by any single-adjoint-spurion mixing of the form $\epsilon^a\,\Psi^\dagger T^a\chi$ --- such a contraction saturates at $j(R)=(\dim R-1)/2<\sqrt{C_2(R)}=\sqrt{j(j+1)}$ --- and the natural candidate Wigner--Eckart resolution (in which $\Psi$ and $\Phi_R$ are realised as the singlet and rep-$R$ sub-blocks of a single parent multiplet $R'\supset\mathbf{1}\oplus R$ of $G\supset\mathrm{SU}(2)_W$) is \emph{ruled out} by an explicit calculation. For the doublet ($R=2$), the matrix element of $T^a$ between $\mathbf{1}$ and $\mathbf{2}$ vanishes identically by the SU(2) selection rule $\mathbf{1}\otimes\mathbf{1/2}=\mathbf{1/2}\oplus\mathbf{3/2}$, since $T^a$ is a $j=1$ tensor and the bilinear has no singlet projection. For the triplet ($R=3$) in the natural SO(4) bidoublet parent, the broken generator $T_A^a$ does connect $\mathbf{1}$ and $\mathbf{3}$, but with $\sum_a|\langle\mathbf{1}|T_A^a|\mathbf{3}\rangle|^2 = 3 \neq C_2(3)=2$. More fundamentally, the seed identity $|M^2_{12}|^2 = C_2(R)\,m_0^4 = -M^2_{22}\,m_0^2$ ties two distinct Wilson coefficients (off-diagonal mixing and diagonal mass) to one another, and no single reduced matrix element can encode such a relation. The $\sqrt{C_2(R)}$ off-diagonal therefore remains an engineered Wilson-coefficient relation; equivalent reformulations include a non-renormalisable operator $S\sqrt{\Phi_R^\dagger T^a T^a\Phi_R}$ or a rep-dependent kinetic normalisation $\tilde\Phi_R\equiv\Phi_R/\sqrt{C_2(R)}$, but neither is a derivation.

\paragraph{(b) $M_W$ is anchored, not predicted.} The ``prediction'' $M_W=80.369$~GeV is a one-parameter fit of $m_0$ to $M_W$. What the model genuinely predicts at $\sim 0.01\%$ accuracy is the \emph{ratio} $M_W:M_Z:m_{h,\mathrm{bare}}$ --- equivalently, the small dispersion of the four eigenvalues~(\ref{eq:eig23}) across the doublet--triplet pair. The Standard Model already predicts $M_W$ from $(M_Z,G_F,\alpha)$ at sub-permille accuracy via electroweak radiative corrections; the Casimir-locked scheme adds no new content at this level, beyond reparametrising one experimental ratio.

\paragraph{(c) $\Delta$ is a counterterm-like tunable.} The spurion $\Delta=(26.42\,\mathrm{GeV})^2$ is exactly the value required to drag $m_{h,\mathrm{bare}}$ onto $m_{h,\mathrm{phys}}$. No symmetry forbids $\Delta=0$, and no first-principles calculation fixes its magnitude inside the seed. The M4 quadratic Coleman--Weinberg estimate (Sec.~\ref{sec:mech}) lands within $1\%$ of the required value but is regulator-dependent; the M3 dim-6 EFT alternative lands at TeV-scale $\Lambda_d$ but the Wilson coefficient is free. Either reading produces $\Delta$ of the right order without uniquely fixing its value.

\paragraph{(d) The triplet lower eigenvalue is unassigned.} The doublet eigenvalues~(\ref{eq:fitB}) are absorbed into $M_W^2$ (upper) and $m_{h,\mathrm{bare}}^2$ (lower). For the triplet, the upper eigenvalue gives $M_Z$ (Eq.~(\ref{eq:MZ})); the lower eigenvalue yields $|\lambda_-(3)|^{1/2}\,m_0 = 176.149$~GeV, with no observable currently assigned to it. The two closest physical neighbours are the Fermi scale $v/\sqrt{2}=174.10$~GeV ($+1.17\%$) and the top-quark pole mass $m_t=172.69$~GeV ($+2.00\%$); since $m_t=y_t\,(v/\sqrt{2})$ with $y_t\simeq 0.99$, these are intrinsically close and the model cannot distinguish them. Either reading would be physically natural, but the residual ($\sim 1$--$2\%$) is two orders of magnitude worse than the $\sim 10^{-4}$ accuracy of the other three eigenvalue slots, so this is at best a suggestive coincidence rather than a fourth precision prediction. The most conservative reading is that the construction overcounts by one eigenvalue and the $176$~GeV slot is genuinely unassigned. Mesonic interpretations are excluded: the heaviest QCD bound state is $\Upsilon(11020)\simeq 11.02$~GeV, more than an order of magnitude below the orphan.

\paragraph{(e) The model is silent on the $M_W\to M_Z$ limit and on the electromagnetic coupling.} The Casimir-locked seesaw is purely algebraic: it produces a mass spectrum from representation invariants $C_2(R)$, but it introduces no gauge fields, no couplings $g,g'$, and no Weinberg angle as a function of couplings. The ratio
\begin{equation}\label{eq:cosw}
\frac{M_W^2}{M_Z^2}\;=\;\frac{\lambda_+(2)}{\lambda_+(3)}\;=\;\frac{(\sqrt{57}-3)/8}{\sqrt{3}-1}\;=\;0.7770
\end{equation}
is \emph{hard-wired} by the doublet and triplet Casimirs; numerically it matches PDG $\cos^2\theta_W$ to $10^{-4}$, but it is fixed by the choice of representations rather than dialled by a continuous parameter. Two consequences follow. (i)~The model admits no continuous limit in which $M_W\to M_Z$: custodial-style degeneracy would require choosing two representations whose upper eigenvalues coincide, which is impossible for distinct $\mathrm{SU}(2)$ irreps because $\lambda_+(R)$ is strictly monotonic in $C_2(R)$. There is no Casimir-side analogue of the SM custodial limit $g'\to 0$. (ii)~The unbroken $\mathrm{U}(1)_{\mathrm{EM}}$ subgroup, the photon, the running of $\alpha_{\mathrm{em}}$, and the matching $e=g\sin\theta_W$ have no analogue inside the seed. They must be supplied by an external SM-like gauge layer wrapped around the algebraic mass-spectrum core.

\paragraph{(f) Terminology.} The ``trace-preserving spurion'' description in earlier formulations is technically incorrect (Sec.~\ref{sec:spurion}). The operative perturbation is rank-1 along the tachyonic eigenvector. The corrected terminology should be used in all subsequent discussion.

\paragraph{(g) The fermion sector is out of scope.} This paper is restricted to the electroweak boson spectrum. A previous version of this work proposed to extend the same algebraic skeleton to charged leptons via a Hermitian $\mathbb{Z}_3$ circulant subject to the Koide constraint $r=a/\sqrt{2}$, and to the down-quark $(s,c,b)$ tuple. Both extensions have been removed because: (i)~the Koide constraint is \emph{imposed} on the circulant as an algebraic identity rather than \emph{derived} from any Lagrangian coupled to the seed of Eq.~(\ref{eq:seed}); (ii)~the lepton ``predictions'' require fixing the centre $a=\tfrac{1}{3}\sum\sqrt{m_i}$ from the full observed triplet plus anchoring the phase to $\sqrt{m_\tau}$, so they are a constrained reconstruction rather than a one-input prediction; (iii)~on the quark side the apparent exactness $Q(m_s,m_c,m_b)=2/3$ relied on curated running-mass benchmarks rather than a stated scheme. The natural next step is a fermion-mass spurion derivation paralleling the rank-1 boson spurion of Sec.~\ref{sec:spurion} and rooted in the same Lagrangian.

\paragraph{(h) Other unaddressed sectors.} The construction does not address neutrino masses, the CKM and PMNS mixing angles, the strong-CP angle, or the QCD sector.

% ---------------------------------------------------------------
\section{Conclusion}\label{sec:conclusion}

We have presented a minimal algebraic ansatz in which the electroweak gauge-boson and bare-Higgs spectrum follows from two $2\times 2$ Casimir-locked seeds labelled by the two lowest non-trivial $\mathrm{SU}(2)_W$ irreducible representations, sharing a single seed scale $m_0$. Fit~B fixes $m_0=106.5702$~GeV by anchoring the doublet upper eigenvalue to $M_W^2$ and predicts $M_Z$ and $m_{h,\mathrm{bare}}$ at the $10^{-4}$ level. A rank-1 eigenstate-aligned spurion $\Delta=(26.4189\,\mathrm{GeV})^2$ --- naturally generated by the one-loop top Coleman--Weinberg shift at cutoff $m_0$, with Haba \emph{et al.}~\cite{Haba2016} supplying the closest mechanistic precedent for such a radiatively-generated bosonic-seesaw spurion, and equivalently encoded by a dim-6 $|H|^4$ EFT at $\Lambda_d\sim 1$--$2$~TeV --- lifts the bare Higgs to the physical pole. The construction is purely electroweak in scope; the fermion sector, neutrino masses, mixing angles, and gauge couplings are not addressed.

The next derivable items are: (i)~a UV mechanism that fixes $b=1$ from a symmetry rather than as a postulate; (ii)~a derivation of Rule~A from a parent group structure such as the $\mathrm{SO}(32)\supset\mathrm{SO}(28)\times\mathrm{SO}(4)$ chain that produces $\{2,3\}$ as the \emph{minimal} doublet/triplet content; (iii)~since the natural Wigner--Eckart derivation of the $\sqrt{C_2(R)}$ off-diagonal in a single parent multiplet is ruled out (Sec.~\ref{sec:caveats}(a)), the question becomes whether \emph{two} operators (mixing and diagonal mass) generated by the same UV portal --- e.g.\ a higher-dimensional operator $\mathcal{O}_{\rm UV}\sim(\Psi^\dagger T^a\chi)(\bar Q T^a Q)/\Lambda_{\rm UV}^2$ shared between the elementary--composite mixing and the composite mass term --- can naturally produce the Wilson-coefficient relation $|m_{12}|^2/|m_{22}| = m_0^2$ characteristic of the seed; this is the central open derivation problem and the natural target for a partial-compositeness UV completion in the spirit of Ref.~\cite{Caracciolo2013PC}; (iv)~a fermion-mass spurion derivation that mirrors the rank-1 $\Delta$ structure of Sec.~\ref{sec:spurion} and ties an extended Yukawa sector to the seed at the Lagrangian level; (v)~a Lagrangian completion that wraps $\mathrm{SU}(2)_W\times\mathrm{U}(1)_Y$ gauge fields around the algebraic mass-spectrum core and recovers $\cos^2\theta_W=0.7770$ as a dynamical output, addressing the silence noted in Sec.~\ref{sec:caveats}(e).

% ---------------------------------------------------------------
\begin{acknowledgments}
The author thanks the anonymous reviewers for constructive comments. This work made substantial use of generative-AI tools (Anthropic Claude, model Opus~4.7) for literature search, symbolic eigenvalue verification with \texttt{sympy} and \texttt{mpmath}, cross-checking of group-theoretic identities, drafting and revision of prose, and exploratory analysis of candidate Lagrangian completions. All physics claims, equations, numerical values, and citations have been reviewed and verified independently by the author. The use of generative AI is disclosed here in line with APS author guidelines.
\end{acknowledgments}

% ---------------------------------------------------------------
\begin{thebibliography}{9}

\bibitem{Rivero2024SO32}
A.~Rivero,
\emph{An interpretation of scalars in SO(32)},
Eur.\ Phys.\ J.\ C \textbf{84}, 1058 (2024),
\href{https://arxiv.org/abs/2407.05397}{arXiv:2407.05397}.

\bibitem{Rivero2005sBoot}
A.~Rivero,
\emph{Supersymmetry with composite bosons},
arXiv:hep-ph/0512065 (2005).
[preprint, no journal publication located on InspireHEP]

\bibitem{Slansky1981}
R.~Slansky,
\emph{Group Theory for Unified Model Building},
Phys.\ Rep.\ \textbf{79}, 1--128 (1981).

\bibitem{PDG2024}
Particle Data Group (R.~L.~Workman \emph{et al.}),
\emph{Review of Particle Physics},
Prog.\ Theor.\ Exp.\ Phys.\ \textbf{2024}, 083C01 (2024).

\bibitem{CalmetOliver2007}
X.~Calmet and J.~F.~Oliver,
\emph{A Seesaw Mechanism in the Higgs Sector},
Eur.\ Phys.\ J.\ C \textbf{50}, 113--118 (2007),
\href{https://arxiv.org/abs/hep-ph/0606209}{arXiv:hep-ph/0606209}.

\bibitem{Kim2005}
H.~D.~Kim,
\emph{Electroweak Symmetry Breaking from SUSY Breaking with Bosonic See-Saw Mechanism},
Phys.\ Rev.\ D \textbf{72}, 055015 (2005),
\href{https://arxiv.org/abs/hep-ph/0501059}{arXiv:hep-ph/0501059}.

\bibitem{Chivukula1999}
R.~S.~Chivukula, B.~A.~Dobrescu, H.~Georgi, and C.~T.~Hill,
\emph{Top Quark Seesaw Theory of Electroweak Symmetry Breaking},
Phys.\ Rev.\ D \textbf{59}, 075003 (1999),
\href{https://arxiv.org/abs/hep-ph/9809470}{arXiv:hep-ph/9809470}.

\bibitem{Haba2016}
N.~Haba, H.~Ishida, R.~Takahashi, and Y.~Yamaguchi,
\emph{Gauge coupling unification in a classically scale invariant model},
Phys.\ Lett.\ B \textbf{754}, 349--352 (2016),
\href{https://arxiv.org/abs/1508.06828}{arXiv:1508.06828}.

\bibitem{Caracciolo2013PC}
F.~Caracciolo, A.~Parolini, and M.~Serone,
\emph{UV Completions of Composite Higgs Models with Partial Compositeness},
J.\ High Energy Phys.\ \textbf{02}, 066 (2013),
\href{https://arxiv.org/abs/1211.7290}{arXiv:1211.7290}.

\end{thebibliography}

\end{document}
